% \Fidel{\Cref{eq:spp-probabilities} and mapping of \(\lambda_{\vec{\alpha}}\) (Pauli probabilities) 
% from \(\tilde{\chi}_{\vec{\xi}}\) can be seen as the multi-time generalisation of Section II B. of~\cite{flammiaEfficientEstimationPauli2020}}
\section{Proof of process-separability}\label{app:proof-pt-to-spp}
Here, we provide the proof of \cref{thm:pt-to-spp}.
\begin{proof}
    The dynamics of a \(k\)-slot process tensor are described by its Choi operator \(\Upsilon_{0:k}\in\mathcal{B}(\mathcal{H}_{\Upsilon_{0:k}})\), where \(\mathcal{H}_{\Upsilon_{0:k}}:=\bigotimes_{j=0}^k(\mathcal{H}_{\inp_j}\otimes\mathcal{H}_{\out_j})\). We expand \(\Upsilon_{0:k}\) in the Pauli basis on each input-output pair as
    \begin{equation}\label{eq:pauli-expansion}
        \Upsilon_{0:k} = \sum_{\vec{\mu},\vec{\nu}} \chi_{\vec{\mu},\vec{\nu}} \bigotimes_{j=0}^k (P_{\mu_j}\otimes P_{\nu_j}),
    \end{equation}
    denoting \(\vec{\mu} = (\mu_0,\mu_1,\cdots,\mu_k)\) and \(\vec{\nu} = (\nu_0,\nu_1,\cdots,\nu_k)\). Here the Pauli operators \(P_{\mu_j}\) and \(P_{\nu_j}\) act on \(\mathcal{H}_{\inp_j}\) and \(\mathcal{H}_{\out_j}\), respectively. Up to a convention-dependent normalisation factor, \(\chi_{\vec{\mu},\vec{\nu}}\) are real coefficients given by
    \begin{equation}
        \chi_{\vec{\mu},\vec{\nu}} = \Tr\left[P_{\vec{\mu},\vec{\nu}}\Upsilon_{0:k}\right], \quad
        P_{\vec{\mu},\vec{\nu}} \coloneq \bigotimes_{j=0}^k (P_{\mu_j}\otimes P_{\nu_j}).
    \end{equation}
    The multi-time Pauli twirl \(\mathcal{T}_P^{(k)}\) is defined in \cref{def:multi-time-twirl}, which we restate here,
    \begin{equation}
        \Upsilon_{0:k}^{\mathcal{T}_P} := \mathcal{T}_P^{(k)}(\Upsilon_{0:k}) = \frac{1}{|\mathbb{P}^{(n)}|^{k+1}} \sum_{Q_0,\dots,Q_k\in\mathbb{P}^{(n)}} \bigg(\bigotimes_{j=0}^k Q_j\otimes Q_j\bigg) \Upsilon_{0:k} \bigg(\bigotimes_{j=0}^k Q_j\otimes Q_j\bigg).
    \end{equation}
    Substituting the Pauli expansion from \cref{eq:pauli-expansion} yields
    \begin{equation}\label{eq:pauli-sub}
        \Upsilon_{0:k}^{\mathcal{T}_P} = \frac{1}{|\mathbb{P}^{(n)}|^{k+1}} \sum_{Q_0,\dots,Q_k\in\mathbb{P}^{(n)}} \sum_{\vec{\mu},\vec{\nu}} \chi_{\vec{\mu},\vec{\nu}} \bigotimes_{j=0}^k (Q_j P_{\mu_j} Q_j \otimes Q_j P_{\nu_j} Q_j).
    \end{equation}
    Using the Pauli conjugation averaging identity,
    \begin{equation}
        \frac{1}{|\mathbb{P}^{(n)}|}\sum_{Q\in\mathbb{P}^{(n)}} (Q P_{\mu} Q) \otimes (Q P_{\nu} Q) = \delta_{\mu,\nu} (P_{\mu} \otimes P_{\nu}),
    \end{equation}
    all cross terms in \(\chi_{\vec{\mu},\vec{\nu}}\) with \(\vec{\mu} \neq \vec{\nu}\) in \cref{eq:pauli-sub} vanish, simplifying the expression to
    \begin{equation}
        \Upsilon_{0:k}^{\mathcal{T}_P} = \sum_{\vec{\xi}} \tilde{\chi}_{\vec{\xi}} \bigotimes_{j=0}^k (P_{\xi_j} \otimes P_{\xi_j}),
    \end{equation}
    where we have relabelled the remaining coefficients as \(\tilde{\chi}_{\vec{\xi}} = \chi_{\vec{\xi},\vec{\xi}}\), \(\vec{\xi} = (\xi_0, \xi_1, \dots, \xi_k)\).

    To prove the absence of quantum temporal correlations, we must show that \(\Upsilon_{0:k}^{\mathcal{T}_P}\) is separable across the bipartitions between distinct time steps. The key insight is that the Bell basis states, \(\ket{\beta_i} \in \{\ket{\Phi^+}, \ket{\Phi^-}, \ket{\Psi^+}, \ket{\Psi^-}\}\) (for a single qubit example), are simultaneous eigenvectors of operators of the form \(P_\xi \otimes P_\xi\). Consequently, \(\Upsilon_{0:k}^{\mathcal{T}_P}\) is diagonal in the Bell product basis \(\ket{\mathcal{B}_{\vec{\alpha}}}=\bigotimes_{j=0}^k \ket{\beta_{\alpha_j}}\). Explicitly, we perform the change of basis
    \begin{equation}
        \Upsilon_{0:k}^{\mathcal{T}_P} = \sum_{{\vec{\alpha}}{\vec{\beta}}} \lambda_{\vec{\alpha},\vec{\beta}} \ketbra{\mathcal{B}_{\vec{\alpha}}}{\mathcal{B}_{\vec{\beta}}},
        \quad
        \lambda_{\vec{\alpha},\vec{\beta}} \coloneq \bra{\mathcal{B}_{\vec{\alpha}}} 
        \sum_{\vec{\xi}} \tilde{\chi}_{\vec{\xi}} \bigotimes_{j=0}^k (P_{\xi_j} \otimes P_{\xi_j})
        \ket{\mathcal{B}_{\vec{\beta}}},
    \end{equation}
    and since \(\ket{\mathcal{B}_{\vec{\alpha}}}\) are eigenvectors of \(\bigotimes_{j=0}^k (P_{\xi_j} \otimes P_{\xi_j})\), we have \(\lambda_{\vec{\alpha},\vec{\beta}} = 0\) for all \(\vec{\alpha} \neq \vec{\beta}\). Hence, we may write
    \begin{equation}\label{eq:bell-basis-decomposition}
        \Upsilon_{0:k}^{\mathcal{T}_P} = \sum_{\vec{\alpha}} \lambda_{\vec{\alpha}} \ketbra{\mathcal{B}_{\vec{\alpha}}}{\mathcal{B}_{\vec{\alpha}}},
        \quad
        \lambda_{\vec{\alpha}} \coloneq \bra{\mathcal{B}_{\vec{\alpha}}} 
        \sum_{\vec{\xi}} \tilde{\chi}_{\vec{\xi}} \bigotimes_{j=0}^k (P_{\xi_j} \otimes P_{\xi_j})
        \ket{\mathcal{B}_{\vec{\alpha}}}.
    \end{equation}
    Because \(\Upsilon_{0:k}^{\mathcal{T}_P}\) is positive semidefinite and the vectors \(\ket{\mathcal{B}_{\vec{\alpha}}}\) form an eigenbasis, the corresponding eigenvalues \(\lambda_{\vec{\alpha}}\) in \cref{eq:bell-basis-decomposition} are non-negative. Moreover, since \(\Upsilon_{0:k}^{\mathcal{T}_P}\) is diagonal in the Bell product basis, its trace is the sum of these eigenvalues,
    \begin{equation}
        \sum_{\vec{\alpha}} \lambda_{\vec{\alpha}} = \Tr[\Upsilon_{0:k}^{\mathcal{T}_P}] = \Tr[\Upsilon_{0:k}],
    \end{equation}
    where the equality follows from the trace preservation of the twirl. Hence, \(\Upsilon_{0:k}^{\mathcal{T}_P}\) is manifestly process-separable (cf. \cref{eq:process-separable}), and the normalised weights \(p_{\vec{\alpha}} = \lambda_{\vec{\alpha}} / \Tr[\Upsilon_{0:k}]\) define a probability distribution as in \cref{eq:spp-probabilities}.

    Finally, the coefficients \(\lambda_{\vec{\alpha}}\) can be related to the Pauli decomposition coefficients \(\tilde{\chi}_{\vec{\xi}}\) via a Walsh-Hadamard transform,
    \begin{equation}
        \lambda_{\vec{\alpha}} = \sum_{\vec{\xi}} \tilde{\chi}_{\vec{\xi}} \prod_{j=0}^k M_{\alpha_j\xi_j}, \quad M_{\alpha\xi} \coloneq \bra{\beta_\alpha}(P_\xi\otimes P_\xi)\ket{\beta_\alpha}.
    \end{equation}
\end{proof}